% Part: normal-modal-logic
% Chapter: completeness
% Section: canonical-models

\documentclass[../../../include/open-logic-section]{subfiles}

\begin{document}

\olfileid[fr]{nml}{com}{cmd}

\olsection{Modèles canoniques}

Pour un système modal cohérent~$\Sigma$, le \emph{modèle canonique} est un modèle
particulier~$\mModel{M^\Sigma}$ dont les mondes sont tous les
ensembles $\Sigma$-cohérents complets. Sa relation
d'accessibilité~$R^\Sigma$ et sa valuation $V^\Sigma$ sont définies de
façon que les !!{formula}s vraies en un monde~$\Delta$ soient
exactement les !!{formula}s qui composent~$\Delta$.

\begin{defn}
  Soit $\Sigma$ une logique modale normale et cohérente. Le \emph{modèle canonique}
  de $\Sigma$ est $\mModel{M}^\Sigma = \tuple{W^\Sigma, R^\Sigma,
    V^\Sigma}$, où :
  \begin{enumerate}
  \item $W^\Sigma = \Setabs{ \Delta }{\Delta \text{
      est $\Sigma$-cohérent complet} }$.
  \item $R^\Sigma \Delta\Delta'$ vaut si et seulement si
    \iftag{prvBox}{$\Box^{-1}\Delta \subseteq
      \Delta'$}{$\Diamond\Delta' \subseteq \Delta$}.
  \item $V^\Sigma(p) = \Setabs{\Delta}{p \in \Delta}$.
  \end{enumerate}
\end{defn}

L'ensemble vide est $\Sigma$-cohérent. La
\olref[lin]{thm:lindenbaum} lui donne une extension $\Sigma$-cohérente
complète. Ainsi, $W^\Sigma$ contient au moins un monde.
\end{document}
